\documentclass[11pt,a4paper]{article}

\usepackage[margin=2cm]{geometry}
\usepackage{enumitem}
\usepackage[hidelinks]{hyperref}
\usepackage[T1]{fontenc}
\usepackage{titlesec}

\setlength{\parindent}{0pt}
\pagenumbering{gobble}

\titleformat{\section}{\large\bfseries}{}{0em}{}[\titlerule]

\begin{document}

% ================= HEADER =================
\begin{center}
{\Large \textbf{ING. MIGUEL ANGEL ANAY GOMEZ}}\\
Lima, Perú\\
\textbf{Senior Frontend Developer | React | Next.js | TypeScript | APIs | Performance | UX/UI}\\
\href{mailto:contacto@miguel-anay.nom.pe}{contacto@miguel-anay.nom.pe} \;|\;
\href{https://www.linkedin.com/in/miguel-anay/}{LinkedIn} \;|\;
+51 936 327 402 \;|\;
\href{https://portfolio.miguel-anay.nom.pe}{Portafolio}
\end{center}

\vspace{0.3cm}

% ================= PERFIL =================
\section*{Perfil Profesional}

Ingeniero de Sistemas titulado con más de \textbf{8 años de experiencia en desarrollo Frontend} y más de \textbf{12 años en desarrollo de software}.  
Especializado en la construcción de aplicaciones web \textbf{mobile-first}, de alto rendimiento y excelente experiencia de usuario utilizando \textbf{React, Next.js y TypeScript}.  
Experiencia sólida en consumo de \textbf{APIs REST}, manejo de estados complejos, optimización de performance, accesibilidad y buenas prácticas de arquitectura frontend.  
Acostumbrado a trabajar bajo metodologías \textbf{Agile/Scrum}, colaborando con equipos de diseño UX/UI, backend y negocio.  
Disponibilidad para trabajo híbrido o presencial en Lima.

% ================= EXPERIENCIA =================
\section*{Experiencia Profesional}

\textbf{Yiwu Import Corporation} \hfill \textit{Abr 2025 – Dic 2025}\\
\textit{Senior Frontend / Full Stack Developer}\\
\textit{Desarrollo de aplicaciones web mobile-first enfocadas en performance y experiencia de usuario}

\begin{itemize}[leftmargin=*]
  \item Desarrollé aplicaciones frontend utilizando \textbf{React} y \textbf{Vue.js} con enfoque mobile-first.
  \item Participé en el desarrollo del sistema de facturación electrónica \textbf{FactuFácil} (\href{https://factufacil.pe/}{factufacil.pe}) utilizando \textbf{Vue.js y Laravel}.
  \item Implementé interfaces basadas en diseños \textbf{UX/UI}, asegurando consistencia visual y usabilidad.
  \item Construí componentes reutilizables, escalables y mantenibles.
  \item Integré el frontend con \textbf{APIs REST} backend.
  \item Implementé y desarrollé \textbf{addons en Odoo} para facturación electrónica e integración con sistemas contables.
  \item Optimicé tiempos de carga y experiencia del usuario.
  \item Colaboré activamente con equipos técnicos y de negocio.
\end{itemize}

\textbf{Logros:} Mejoré la experiencia de usuario y la escalabilidad del frontend en plataformas de negocio.

\vspace{6pt}

\textbf{Encora} \hfill \textit{Nov 2020 – Abr 2025}\\
\textit{Frontend / Full Stack Developer – Cliente Profuturo AFP}\\
\textit{Desarrollo de aplicaciones web financieras de alta disponibilidad}

\begin{itemize}[leftmargin=*]
  \item Desarrollé aplicaciones web con \textbf{React y Angular}, construyendo componentes reutilizables y escalables bajo metodología \textbf{Agile/Scrum}.
  \item Implementé flujos con manejo de estado, validaciones y control de permisos, integrando \textbf{APIs REST} seguras en plataformas financieras críticas.
  \item Optimicé performance, accesibilidad y experiencia de usuario en sistemas de alta disponibilidad.
  \item Contribuí al desarrollo de \textbf{Agencia Virtual}, optimizando trámites digitales y reduciendo la carga operativa presencial.
  \item Participé en la implementación de \textbf{Clave Web}, fortaleciendo la seguridad digital y reduciendo incidencias de soporte en más de un \textbf{40\%}.
  \item Desarrollé funcionalidades en \textbf{Disfruta Profuturo}, incrementando el engagement digital de afiliados.
\end{itemize}

\textbf{Logros:} Reduje incidencias y mejoré la experiencia del usuario, disminuyendo consultas operativas en más del \textbf{40\%}.

\vspace{6pt}

\textbf{Prosegur} \hfill \textit{Abr 2018 – Mar 2020}\\
\textit{Frontend / Full Stack Developer}

\begin{itemize}[leftmargin=*]
  \item Desarrollé aplicaciones web utilizando \textbf{React.js} y \textbf{PHP Laravel}, integradas con \textbf{MySQL y SQL Server}.
  \item Implementé soluciones de monitoreo y trazabilidad operativa con integración GPS y mapas interactivos.
  \item Desarrollé \textbf{Web Audit}, mejorando control interno y reduciendo incidentes operativos.
  \item Contribuí en \textbf{Web PER}, sistema de monitoreo en tiempo real con alertas automáticas.
  \item Implementé la \textbf{App de Supervisiones} en \textbf{Android Studio}, fortaleciendo el control operativo en campo.
\end{itemize}

\textbf{Logros:} Digitalicé procesos operativos mejorando el control y reduciendo riesgos.

\vspace{6pt}

\textbf{Medios Industriales} \hfill \textit{Abr 2014 – Mar 2018}\\
\textit{Frontend / Full Stack Developer}

\begin{itemize}[leftmargin=*]
  \item Desarrollé aplicaciones web con \textbf{React.js} y \textbf{PHP Laravel}, integradas con \textbf{MySQL y SQL Server}.
  \item Implementé el sistema \textbf{Control de Tiempos Operativos (CTO)}, optimizando jornadas y productividad.
  \item Desarrollé la plataforma \textbf{Balizas}, sistema de monitoreo GPS con alertas automáticas.
\end{itemize}

\textbf{Logros:} Mejoré el control operativo reduciendo desviaciones en un \textbf{20\%}.

% ================= COMPETENCIAS =================
\section*{Competencias Técnicas}

\textbf{Frontend:} React, Next.js, Vue.js, TypeScript, JavaScript (ES6+)\\
\textbf{Backend:} Laravel, Node.js, APIs REST\\
\textbf{ERP & Integraciones:} Odoo (addons), Facturación electrónica\\
\textbf{Estilos & UI:} Tailwind CSS, Responsive Design, Mobile-first\\
\textbf{Arquitectura:} Clean Architecture, Arquitectura Hexagonal\\
\textbf{Performance:} Lazy Loading, Code Splitting, Core Web Vitals\\
\textbf{Bases de Datos:} SQL Server, MySQL, MongoDB\\
\textbf{DevOps:} Docker, Kubernetes, Azure AZ-900\\
\textbf{Versionamiento:} Git, Azure DevOps\\
\textbf{Metodologías:} Agile / Scrum

% ================= EDUCACION =================
\section*{Educación}

\textbf{Universidad Nacional de Ingeniería (UNI)}\\
Ingeniero Industrial \hfill \textit{2004 – 2012}\\
CIP: 301084

\vspace{4pt}

\textbf{Universidad Privada del Norte}\\
Diplomado en Gerencia de TI\\
\textit{Ene 2025 – Sep 2025}

% ================= IDIOMAS =================
\section*{Idiomas}

Español: Nativo\\
Inglés: Intermedio (lectura técnica)

\end{document}
